% !TEX TS-program = pdflatex
\documentclass[10pt,a4paper]{altacv}

% Marges équilibrées pour 1 page complète
\geometry{left=1.2cm,right=1.2cm,top=1.0cm,bottom=0.8cm,footskip=0pt}

\usepackage[T1]{fontenc}
\usepackage[utf8]{inputenc}
\usepackage[default]{lato}
\usepackage{microtype}
\DisableLigatures{encoding=*, family=*} % ATS-fix: désactive les ligatures en pdfLaTeX

\definecolor{Navy}{HTML}{0A2540}
\definecolor{SlateGrey}{HTML}{2E2E2E}
\definecolor{LightGrey}{HTML}{555555}
\colorlet{heading}{Navy}
\colorlet{accent}{Navy}
\colorlet{emphasis}{SlateGrey}
\colorlet{body}{LightGrey}

\renewcommand{\itemmarker}{{\small\textbullet}}
\renewcommand{\ratingmarker}{\faCircle}

\usepackage[hidelinks]{hyperref}
\hypersetup{
  pdftitle={CV -- Achraf Izelmadene -- Candidature SolarCare},
  pdfauthor={Achraf Izelmadene},
  pdfsubject={Stage PFE Maintenance et Exploitation Centrales Photovoltaïques BtoB},
  pdfkeywords={SolarCare, photovoltaïque, PVsyst, ENERGYSOFT, maintenance PV, Power BI, macros Excel, HTA, BT, stage PFE}
}

% Photo avec diamètre 1.9cm et cadrage naturel
\photo{1.9cm}{assets/imageee.png}

% Style compact pour les sections
\renewcommand{\cvsection}[2][]{%
  \vspace{3pt}%
  {\color{heading}\Large\bfseries #2 \hspace{0.2em}\leaders\vrule height 2.5pt depth -1.25pt\hfill\mbox{}}\par\vspace{2.5pt}
}

\renewcommand{\cvevent}[4]{%
  {\normalsize\bfseries\color{accent}#1\par}\vspace{1pt}
  \ifstrequal{#2}{}{}{%
    {\small\itshape\color{emphasis}#2\ifstrequal{#4}{}{}{ - #4}}%
  }%
  ~\hfill
  \ifstrequal{#3}{}{}{{\small\color{emphasis}#3}}%
  \par\vspace{1.5pt}
}

% Header fusionné avec À Propos
\makeatletter
\newcommand{\target}[1]{\def\@target{#1}}
\newcommand{\apropos}[1]{\def\@apropos{#1}}

\renewcommand{\makecvheader}{%
  \begingroup
    \ifdef{\@photodiameter}{%
    \begin{minipage}[c]{\@photodiameter}
    \if@normalphoto
      \includegraphics[width=\linewidth]{\@photo}
    \else
      \begin{tikzpicture}
        \clip (0,0) circle (0.5\linewidth);
        \node at (0,-0.04\linewidth) {\includegraphics[width=1.15\linewidth]{\@photo}};
      \end{tikzpicture}
    \fi%
    \end{minipage}\hspace{1.2em}%
    \begin{minipage}[c]{\dimexpr\linewidth-\@photodiameter-1.5em}}{}%
    \raggedright\color{emphasis}%
    {\LARGE\bfseries\MakeUppercase{\@name}\par}\vspace{1.5pt}
    {\normalsize\bfseries\color{accent}\@tagline\par}\vspace{3pt}
    \ifdef{\@target}{{\small\textbf{\@target}\par}\vspace{2pt}}{}%
    \ifdef{\@apropos}{{\small\color{emphasis}\@apropos\par}\vspace{3pt}}{}%
    \vspace{2.5pt}%
    {\footnotesize\color{emphasis}\@personalinfo\par}
    \ifdef{\@photodiameter}{%
    \end{minipage}\par}{}%
  \endgroup\vspace{1pt}
}
\makeatother

\begin{document}

\name{Izelmadene Achraf}
\tagline{Élève Ingénieur en Génie Électrique et Systèmes Intelligents (GESI)}

% Profil adapté à l'offre SolarCare (Maintenance & Exploitation PV)
\target{À la recherche d’un stage PFE de 5 à 6 mois dès février 2027 en maintenance, exploitation et rénovation de centrales PV (SolarCare).}
\apropos{Élève ingénieur en génie électrique doté d’une double culture bureau d’études PV (intégration PV sous PVsyst, calculs HTA/BT) et maintenance industrielle (diagnostic, indicateurs). Rigoureux et force de proposition en automatisation d'outils (macros Excel, Power BI, HTML/JS), je souhaite contribuer activement à la performance du parc PV BtoB de SolarCare.}

\personalinfo{%
  \href{mailto:achraf.izelmadene@gmail.com}{achraf.izelmadene@gmail.com}%
  \hspace{0.9em}\textcolor{body!60}{|}\hspace{0.9em}%
  +212699098736%
  \hspace{0.9em}\textcolor{body!60}{|}\hspace{0.9em}%
  Marrakech, Maroc%
  \hspace{0.9em}\textcolor{body!60}{|}\hspace{0.9em}%
  Permis B%
}

\makecvheader

\cvsection{FORMATION}
\cvevent{Étudiant en cycle d’ingénieur en Génie Électrique et Systèmes Intelligents}{École Marocaine des Sciences de l’Ingénieur - Marrakech}{2025 - Présent}{}
{\footnotesize\color{emphasis}Énergies renouvelables : production \& stockage, Conception des installations électriques HTA/BT, Réseaux électriques \& Sécurité, Audit \& Efficacité énergétique, Supervision industrielle \& Digital Twin, Électronique de puissance, Lean Management, Contrôle qualité, Hygiène, Sécurité \& Environnement (HSE), Gestion de projet...}\par\vspace{2pt}

\cvevent{Licence SMP, option Automatique}{Faculté des Sciences Semlalia - Marrakech}{2020 - 2024}{}
{\footnotesize\color{emphasis}Électricité, Électronique analogique, Asservissements linéaires, Automatismes et capteurs, Traitement du signal...}\par

\cvsection{EXPÉRIENCE}
\cvevent{Stagiaire Bureau d’Études}{Core Engineering Solutions}{juil. - août 2026}{Berrchid, Casablanca}
\begin{itemize}[leftmargin=*, label=\textbullet, nosep, itemsep=1.5pt]
  \item Étude d’intégration PV en autoconsommation (PVsyst) : productible, calcul des pertes et schémas unifilaires.
  \item Dimensionnement d’un poste de livraison HTA/BT, bilan de puissance et notes de calcul (Caneco BT, AutoCAD).
\end{itemize}
\vspace{1.5pt}
\cvevent{Stagiaire en Unité de Production}{Menara Prefa}{juil. 2024}{Marrakech}
\begin{itemize}[leftmargin=*, label=\textbullet, nosep, itemsep=1.5pt]
  \item Maintenance préventive et corrective sur lignes de production, diagnostic électrique et suivi des indicateurs.
  \item Migration d’un automate vers Siemens S7 (TIA Portal) et mise à jour de la documentation technique.
\end{itemize}

\cvsection{PROJETS}
{\normalsize\bfseries\color{accent} Détection d’anomalies Smartgrid \& Dashboard (Clustering \& Séries Temporelles)}\par\vspace{1.5pt}
\begin{itemize}[leftmargin=*, label=\textbullet, nosep, itemsep=1.5pt]
  \item Analyse et partitionnement de profils de consommation électrique (KMeans, DBSCAN, GMM) sur données IoT.
  \item Développement d’un dashboard interactif (Python, Streamlit) pour le monitoring et la détection d’anomalies.
\end{itemize}
\vspace{1.5pt}
{\normalsize\bfseries\color{accent} Maintenance Prédictive 4.0 par AutoML (TPOT vs. Auto-WEKA)}\par\vspace{1.5pt}
\begin{itemize}[leftmargin=*, label=\textbullet, nosep, itemsep=1.5pt]
  \item Détection anticipée de défaillances sur données IoT multi-capteurs industrielles (vibrations, température, pression).
  \item Benchmark d’algorithmes (F1-score, Recall) et analyse CND d’huile pour la fiabilisation d’équipements.
\end{itemize}
\vspace{1.5pt}
{\normalsize\bfseries\color{accent} Smart Agriculture IoT \& Dashboard de Supervision (1\textsuperscript{er} Prix Hack The Intelligence)}\par\vspace{1.5pt}
\begin{itemize}[leftmargin=*, label=\textbullet, nosep, itemsep=1.5pt]
  \item Système connecté d’acquisition multi-capteurs (ESP32/LoRa) et plateforme cloud de supervision en temps réel.
  \item Gestion de projet agile en équipe restreinte et soutenance devant un jury d'experts industriels.
\end{itemize}

\cvsection{COMPÉTENCES}
\begin{tabular}{@{}ll@{}}
\textbf{Solaire PV \& BE :} & PVsyst, Caneco BT, AutoCAD, Générateur PV, Onduleurs, Postes HTA/BT \\[1.5pt]
\textbf{Maintenance \& Supervision :} & Supervision PV (Energysoft, SCADA), Diagnostic, Maintenance prév./corr. \\[1.5pt]
\textbf{Outils Métiers \& Data :} & Macros Excel (VBA), Microsoft Power BI, HTML / JavaScript, Python \\[1.5pt]
\textbf{Transversal \& PME 360° :} & Gestion de projet (Dassault Systèmes), Rigueur, Autonomie, Relation client BtoB
\end{tabular}

\cvsection{CERTIFICATIONS}
\begin{itemize}[leftmargin=*, label=\textbullet, nosep, itemsep=1.5pt]
  \item \textbf{Certified 3DEXPERIENCE Collaborative Industry Innovator}~\hfill {\color{emphasis}Dassault Systèmes}
  \item \textbf{Foundations of IoT Systems and Industrial Automation}~\hfill {\color{emphasis}LearnQuest (Coursera)}
  \item \textbf{Sensors and Sensor Circuit Design}~\hfill {\color{emphasis}University of Colorado Boulder (Coursera)}
  \item \textbf{Machine Learning in Production}~\hfill {\color{emphasis}DeepLearning.AI (Coursera)}
\end{itemize}

\cvsection{ACTIVITÉS \& ENGAGEMENT}
\begin{itemize}[leftmargin=*, label=\textbullet, nosep, itemsep=1.5pt]
  \item \textbf{Hack The Intelligence (AISEC 2026) - 1\textsuperscript{er} Prix :} Développement d’un prototype Smart Agriculture IoT.
  \item \textbf{Vision AI Competition (EMSI 2026) - 2\textsuperscript{ème} Prix :} Conception d’un système de traduction en langue des signes (ASL).
  \item \textbf{Congrès COCIA’26 - Présentation Scientifique :} Conception du prototype d’éclairage piézoélectrique SmartStruct.
  \item \textbf{Engagement Associatif \& Bénévolat (2025) :} Soutien scolaire et encadrement pédagogique d'enfants orphelins.
\end{itemize}

\cvsection{LANGUES}
\textbf{Arabe :} Langue maternelle~\hfill \textbf{Français :} DELF B2~\hfill \textbf{Anglais :} Professionnel

\end{document}
